\documentclass[11pt,a4paper]{article}
\usepackage[margin=25mm]{geometry}
\usepackage{amsmath,amssymb,amsthm,hyperref}
\setlength{\emergencystretch}{2em}
\newtheorem{theorem}{Theorem}
\newcommand{\gen}[1]{\langle#1\rangle}
\title{Localized cuts with unselected colours}
\author{Human Compression project: L044 continuation}
\date{}
\begin{document}\maketitle
Let $G=\gen{\rho_0,\ldots,\rho_{r-1}}$ be transitive on $\Omega$,
with involutory generators and string commutations. Eligibility always
means crossing orbits of the full ambient omission $G_j$.
All orbits below include fixed singletons.

\begin{theorem}[Selected colours in a larger interval]\label{larger}
Let $I$ be a contiguous interval, $H=G_I$, and $F\subseteq I$.
Let $V$ be a proper $H$-invariant subset of $\Omega$.
Suppose each colour in $F$ has at least two ambient-eligible edges,
and every ambient-eligible edge of these colours lies in $V$.
Then
\[
                         2|F|\le |V|-o_H(V).             \tag{1}
\]
The statement also holds when $F$ is empty.
\end{theorem}
\begin{proof}
For empty $F$ the inequality is immediate. Otherwise retain the complete
paid paired-cut normal form, selected-colour adapter and ambient-interval
proof, and select two eligible edges of each colour in $F$ on $V$.
We specify the change in their proof rather than importing local
omissions. Long-cycle reduction and ordinary square slides use only
selected colours and the ambient cuts. Their middle move retains a
square of colours $i-1,i+1$ and transports all other selected edges by
$\rho_i$. This is valid even if $i\notin F$, as proved in the
selected-colour adapter. Since the square colours lie in $I$, its
middle $i$ lies in $I$; the move therefore preserves actual $H$-orbits.
All exchanges stay in $V$ by the confinement of every eligible edge.
We obtain a minimum-square selection with at most one square in each
selected component.

Suppose an actual nontrivial $H$-orbit $O\subseteq V$ had no tree
component. Its selected edge count would be $|O|$. The complete
cycle-colour localisation, full-neighbour and partition arguments of
the ambient-interval proof apply unchanged: they need that a square's
middle permutation belongs to $H$, not that it is selected. In the
full-neighbour step an unselected middle is permitted by the same
selected-colour adapter. Hence no full edge from a square vertex leaves
its selected component.

As $O$ is proper, minimize a connector from a selected square through
selected edges in $O$ followed by one full edge leaving $O$, over these
minimum-square selections. Ordinary slides shorten it, including the
accepted repeated-colour case. In the remaining case its first colour
is the square's middle $i$. That colour is selected because it is on
the connector. The exit colour $x$ is outside $I$, since every generator
of $H$ preserves $O$. Both $i-1,i+1$ are in $I$, so $|i-x|\ge2$.
The global middle move transports the exit to an actual exit and
shortens the connector, exactly as in the paid interval proof.
This contradiction supplies a tree in $O$. Fixed vertices in $V$ are
isolated. Each actual $H$-orbit therefore contributes at least one
vertex minus selected edge; summing proves (1).
\end{proof}

\begin{theorem}[A colour set without a one-colour hole]\label{holes}
Let $F$ satisfy the condition
\[
              i-1,i+1\in F\quad\Longrightarrow\quad i\in F. \tag{2}
\]
Put $T=G_F$. If $V$ is proper and $T$-invariant, every colour in $F$
has at least two ambient-eligible edges, and all such edges lie in
$V$, then
\[
                         2|F|\le |V|-o_T(V).             \tag{3}
\]
In particular (2) holds for two intervals separated by at least two
missing colours.
\end{theorem}
\begin{proof}
Run the same proof with actual $T$-orbits. Every square middle belongs
to $F$ by (2), so all middle moves preserve these orbits. A full edge
leaving a $T$-orbit has colour $x\notin F$. If $i$ is a square's
middle, $i-1,i,i+1$ all belong to $F$, so $|i-x|\ge2$ even when
$F$ is disconnected. This is exactly the commutation needed for the
last exit move. All other steps and the singleton accounting are as
in Theorem~\ref{larger}.
\end{proof}

\paragraph{Attribution and scope.}
These are precise extensions of the accepted selected-colour and
ambient-interval proofs supporting Cameron--Fernandes--Leemans--Mixer,
\emph{Highest rank of a polytope for $A_n$}, arXiv:1605.09173v1,
Propositions 5.27 and 5.31. Their full normal-form and equality arguments
remain paid, including the root collaborator's global middle move and
its distinct review. The new work identifies the two index-set
conditions under which orbit preservation and the exit commutation
survive. It assumes neither projected IP nor parity, and does not
identify local and ambient eligibility. No computation, formal theorem
or whole-source credit is claimed.
\end{document}
